\documentclass[a4paper,11pt]{article}

% --- Packages ---
\usepackage[utf8]{inputenc}
\usepackage[T1]{fontenc}
\usepackage[french]{babel}
\usepackage{geometry}
\geometry{margin=2.5cm}
\usepackage{graphicx}
\usepackage{xcolor}
\usepackage{hyperref}
\usepackage{booktabs}
\usepackage{fancyhdr}
\usepackage{listings}
\usepackage{tikz}
\usetikzlibrary{shapes.geometric, arrows.meta, positioning, fit, calc, backgrounds}

% --- Colors ---
\definecolor{primary}{HTML}{2563EB}
\definecolor{secondary}{HTML}{475569}
\definecolor{accent}{HTML}{059669}
\definecolor{lightgray}{HTML}{F1F5F9}
\definecolor{codebg}{HTML}{F8FAFC}

\hypersetup{
    colorlinks=true,
    linkcolor=primary,
    urlcolor=primary,
    citecolor=primary,
}

% --- Header/Footer ---
\pagestyle{fancy}
\fancyhf{}
\fancyhead[L]{\small\textcolor{secondary}{DATA713 -- MLOps}}
\fancyhead[R]{\small\textcolor{secondary}{Telecom Paris}}
\fancyfoot[C]{\thepage}
\renewcommand{\headrulewidth}{0.4pt}

% --- Listings ---
\lstset{
    backgroundcolor=\color{codebg},
    basicstyle=\ttfamily\small,
    breaklines=true,
    frame=single,
    rulecolor=\color{lightgray},
    columns=flexible,
}

% --- TikZ styles ---
\tikzset{
    service/.style={
        rectangle, rounded corners=4pt, minimum width=3.2cm, minimum height=1cm,
        draw=#1!60, fill=#1!8, text=black, font=\small\bfseries,
        align=center, line width=0.8pt
    },
    service/.default=primary,
    storage/.style={
        cylinder, shape border rotate=90, aspect=0.25,
        minimum width=2.4cm, minimum height=1.2cm,
        draw=#1!60, fill=#1!8, text=black, font=\small\bfseries,
        align=center, line width=0.8pt
    },
    storage/.default=secondary,
    arrow/.style={-{Stealth[length=6pt]}, line width=0.8pt, color=secondary},
    dashed arrow/.style={-{Stealth[length=6pt]}, line width=0.8pt, dashed, color=secondary!60},
    group/.style={draw=secondary!30, fill=lightgray, rounded corners=6pt, inner sep=10pt},
}

\begin{document}

% =====================================================================
% TITLE PAGE
% =====================================================================
\begin{titlepage}
\centering
\vspace*{2cm}

{\Huge\bfseries\textcolor{primary}{House Price Prediction}}\\[0.5cm]
{\LARGE\textcolor{secondary}{Plateforme MLOps End-to-End}}\\[2cm]

{\Large DATA713 -- Mastere Specialise IA Expert Data}\\[0.3cm]
{\Large Telecom Paris}\\[2cm]

\rule{0.6\textwidth}{0.4pt}\\[1.5cm]

{\large\textbf{Equipe :} [Noms des membres]}\\[0.5cm]
{\large\textbf{Date :} \today}\\[0.5cm]
{\large\textbf{Branche :} \texttt{Cyprien\_Minio}}\\[2cm]

\vfill
{\small Projet POC en une journee -- King County House Sales Dataset}
\end{titlepage}

% =====================================================================
% TABLE OF CONTENTS
% =====================================================================
\tableofcontents
\newpage

% =====================================================================
% 1. INTRODUCTION
% =====================================================================
\section{Introduction}

Ce rapport presente la conception et la mise en oeuvre d'une plateforme MLOps complete
pour la prediction des prix immobiliers dans le comte de King (Washington, USA).
Le projet couvre l'ensemble du cycle de vie d'un modele de Machine Learning :
de l'acquisition des donnees a la mise en production, en passant par l'entrainement
automatise, le suivi des experiences, le monitoring et le deploiement.

\subsection{Objectifs du projet}

\begin{enumerate}
    \item Extraction et preprocessing des donnees (Kaggle $\rightarrow$ MinIO)
    \item Construction d'un modele ML (XGBoost + Optuna)
    \item Stockage du modele sur un registry (MLflow Model Registry)
    \item Pipeline de reentrainement automatise (Apache Airflow)
    \item Suivi des experiences (MLflow Tracking)
    \item API de serving (FastAPI)
    \item Interface web (Gradio)
    \item Continuous Training avec comparaison de modeles
    \item Dockerisation + Kubernetes + CI/CD
    \item Versionning et documentation (Git + GitHub)
\end{enumerate}

\subsection{Dataset}

\begin{table}[h]
\centering
\begin{tabular}{ll}
\toprule
\textbf{Propriete} & \textbf{Valeur} \\
\midrule
Source & King County House Sales (Kaggle) \\
Taille & 21 613 transactions \\
Features & 21 variables \\
Periode & Mai 2014 -- Mai 2015 \\
Variable cible & \texttt{price} (USD) \\
\bottomrule
\end{tabular}
\caption{Caracteristiques du dataset}
\end{table}

% =====================================================================
% 2. ARCHITECTURE
% =====================================================================
\section{Architecture Globale}

\subsection{Diagramme d'architecture}

\begin{figure}[h]
\centering
\begin{tikzpicture}[node distance=1.2cm and 1.8cm]

    % --- User-facing layer ---
    \node[service=accent] (webapp) {Gradio\\WebApp\\{\scriptsize :7860}};
    \node[service=primary, below=of webapp] (api) {FastAPI\\API\\{\scriptsize :8000}};

    % --- ML layer ---
    \node[service=primary, below left=1.5cm and 2.5cm of api] (mlflow) {MLflow\\Tracking\\{\scriptsize :5000}};
    \node[service=accent, below right=1.5cm and 2.5cm of api] (airflow) {Airflow\\Orchestration\\{\scriptsize :8080}};

    % --- Storage layer ---
    \node[storage=secondary, below=1.5cm of mlflow] (postgres) {PostgreSQL\\{\scriptsize :5432}};
    \node[storage=accent, below=1.5cm of airflow] (minio) {MinIO S3\\{\scriptsize :9000}};

    % --- Monitoring layer ---
    \node[service=orange, right=2.5cm of api] (prometheus) {Prometheus\\{\scriptsize :9090}};
    \node[service=orange, below=of prometheus] (grafana) {Grafana\\{\scriptsize :3000}};

    % --- Arrows ---
    \draw[arrow] (webapp) -- node[right, font=\scriptsize] {HTTP} (api);
    \draw[arrow] (api) -- (mlflow);
    \draw[arrow] (api.east) -- (prometheus.west);
    \draw[dashed arrow] (prometheus) -- (grafana);

    \draw[arrow] (mlflow) -- (postgres);
    \draw[arrow] (airflow) -- (minio);
    \draw[dashed arrow] (airflow.west) -- node[above, font=\scriptsize] {trigger} (mlflow.east);

    % --- Group boxes ---
    \begin{scope}[on background layer]
        \node[group, fit=(webapp)(api), label={[font=\small\bfseries\color{secondary}]above:Serving}] {};
        \node[group, fit=(mlflow)(airflow), label={[font=\small\bfseries\color{secondary}]above:ML Platform}] {};
        \node[group, fit=(postgres)(minio), label={[font=\small\bfseries\color{secondary}]below:Stockage}] {};
        \node[group, fit=(prometheus)(grafana), label={[font=\small\bfseries\color{secondary}]above:Monitoring}] {};
    \end{scope}

\end{tikzpicture}
\caption{Architecture de la plateforme MLOps}
\label{fig:architecture}
\end{figure}

\subsection{Justification des choix technologiques}

\begin{table}[h]
\centering
\small
\begin{tabular}{p{3cm}p{3cm}p{7cm}}
\toprule
\textbf{Composant} & \textbf{Choix} & \textbf{Justification} \\
\midrule
Modele ML & XGBoost & Performant sur donnees tabulaires, rapide a entrainer, interpretable via feature importance. Largement utilise en production. \\
\addlinespace
Hyperparametres & Optuna & Optimisation bayesienne efficace, pruning automatique des essais non prometteurs, integration native avec XGBoost. \\
\addlinespace
Tracking & MLflow & Standard de facto pour le suivi d'experiences ML. Model Registry integre pour la gestion des versions de modeles. \\
\addlinespace
Object Storage & MinIO & Compatible S3, leger, auto-heberge. Permet de decoupler le stockage des donnees du systeme de fichiers local. \\
\addlinespace
Orchestration & Airflow & DAGs Python natifs, interface web, scheduling flexible. Adapte pour les pipelines de reentrainement periodiques. \\
\addlinespace
API & FastAPI & Haute performance (async), validation automatique (Pydantic), documentation OpenAPI generee, support natif de Prometheus. \\
\addlinespace
WebApp & Gradio & Prototypage rapide d'interfaces ML, composants preconstruits, ideal pour une demo/POC. \\
\addlinespace
Base de donnees & PostgreSQL & Backend robuste partage entre MLflow (tracking store) et Airflow (metadata store). \\
\addlinespace
Monitoring & Prometheus + Grafana & Stack standard pour le monitoring d'applications. Prometheus scrape les metriques, Grafana les visualise. \\
\addlinespace
CI/CD & GitHub Actions & Integration native avec GitHub, gratuit pour les repos publics, workflows YAML declaratifs. \\
\addlinespace
Conteneurisation & Docker Compose + K8s & Compose pour le dev local (simplicite), Kubernetes pour la production (scalabilite, resilience). \\
\bottomrule
\end{tabular}
\caption{Justification des choix technologiques}
\end{table}

% =====================================================================
% 3. PIPELINE DE DONNEES
% =====================================================================
\section{Pipeline de Donnees}

\subsection{Acquisition}

Les donnees sont telechargees depuis Kaggle via l'API officielle (\texttt{src/data/download.py})
puis stockees sur MinIO dans le bucket \texttt{mlops-data}. L'authentification Kaggle
se fait via les variables d'environnement \texttt{KAGGLE\_USERNAME} et \texttt{KAGGLE\_KEY}.

\subsection{Preprocessing}

Le preprocessing (\texttt{src/data/preprocess.py}) suit un pipeline en 5 etapes :

\begin{enumerate}
    \item \textbf{Nettoyage} -- Parsing des dates, suppression de la colonne \texttt{id},
          suppression de l'outlier a 33 chambres, imputation des \texttt{NaN} pour \texttt{waterfront}.
    \item \textbf{Feature engineering} -- Creation de variables derivees :
          \texttt{house\_age}, \texttt{is\_renovated}, \texttt{years\_since\_renovation},
          \texttt{has\_basement}, ratios de surface (habitation vs voisins, terrain vs voisins).
    \item \textbf{Split train/test} -- 80/20, effectue \emph{avant} l'encodage pour eviter le data leakage.
    \item \textbf{Target Encoding} -- Encodage du code postal via \texttt{TargetEncoder}
          (ajuste sur le train uniquement).
    \item \textbf{Transformation cible} -- \texttt{log1p(price)} pour l'entrainement,
          \texttt{expm1()} a l'inference.
\end{enumerate}

Les artefacts de preprocessing (encodeur, noms de features) sont sauvegardes sur MinIO
pour etre reutilises a l'inference.

% =====================================================================
% 4. MODELE
% =====================================================================
\section{Modele et Entrainement}

\subsection{XGBoost + Optuna}

L'entrainement (\texttt{src/model/train.py}) utilise XGBoost avec optimisation
bayesienne des hyperparametres via Optuna (10 essais par defaut).

Hyperparametres optimises :
\begin{itemize}
    \item \texttt{max\_depth} : 3--12
    \item \texttt{learning\_rate} : 0.01--0.3
    \item \texttt{n\_estimators} : 100--1000
    \item \texttt{subsample} : 0.6--1.0
    \item \texttt{colsample\_bytree} : 0.6--1.0
    \item \texttt{reg\_alpha} / \texttt{reg\_lambda} : regularisation L1/L2
\end{itemize}

\subsection{Metriques d'evaluation}

\begin{table}[h]
\centering
\begin{tabular}{ll}
\toprule
\textbf{Metrique} & \textbf{Description} \\
\midrule
RMSE (prix reel) & Erreur quadratique moyenne sur les prix en dollars \\
MAE (prix reel) & Erreur absolue moyenne \\
R\textsuperscript{2} & Coefficient de determination \\
MAPE & Pourcentage d'erreur moyen \\
\bottomrule
\end{tabular}
\caption{Metriques de performance du modele}
\end{table}

\subsection{MLflow Tracking}

Chaque run d'entrainement enregistre dans MLflow :
\begin{itemize}
    \item Les hyperparametres (params)
    \item Les metriques sur le jeu de test
    \item Le modele serialise (artefact XGBoost)
    \item Les metadonnees (nombre d'echantillons, duree)
\end{itemize}

Le modele est enregistre dans le \textbf{Model Registry} sous le nom
\texttt{house-price-xgboost}.

% =====================================================================
% 5. CONTINUOUS TRAINING
% =====================================================================
\section{Continuous Training}

\subsection{Orchestration Airflow}

Deux DAGs Airflow gerent le cycle de vie des donnees et du modele :

\begin{table}[h]
\centering
\begin{tabular}{lll}
\toprule
\textbf{DAG} & \textbf{Schedule} & \textbf{Role} \\
\midrule
\texttt{data\_pipeline} & Hebdomadaire & Verifier fraicheur des donnees, telecharger, valider \\
\texttt{training\_pipeline} & Declenche & Entrainer, comparer, promouvoir, alerter \\
\bottomrule
\end{tabular}
\caption{DAGs Airflow}
\end{table}

\subsection{Comparaison de modeles}

Le pipeline de training implemente une logique de \textbf{promotion conditionnelle} :

\begin{enumerate}
    \item Le nouveau modele est entraine et ses metriques sont calculees sur le jeu de test.
    \item Le modele en production actuel est charge depuis le Model Registry MLflow.
    \item Les deux modeles sont compares sur le RMSE (prix reels).
    \item Le nouveau modele est promu en production \textbf{uniquement s'il est meilleur}.
    \item Une alerte email est envoyee avec le resultat (promotion ou rejet).
\end{enumerate}

Cette approche garantit une \textbf{non-regression} : le modele en production
ne se degrade jamais suite a un reentrainement automatique.

% =====================================================================
% 6. SERVING
% =====================================================================
\section{Serving}

\subsection{API FastAPI}

L'API (\texttt{src/api/app.py}) expose :
\begin{itemize}
    \item \texttt{POST /predict} -- Prediction de prix a partir des features
    \item \texttt{GET /health} -- Healthcheck (etat du modele, URI MLflow)
    \item \texttt{GET /metrics} -- Metriques Prometheus
    \item \texttt{GET /docs} -- Documentation OpenAPI (Swagger)
\end{itemize}

Au demarrage, l'API charge le modele depuis MLflow et les artefacts de preprocessing
(encodeur, noms de features) depuis MinIO.

\subsection{Interface Gradio}

L'interface web (\texttt{src/webapp/app.py}) permet aux utilisateurs de saisir
les caracteristiques d'un bien immobilier et d'obtenir une estimation de prix.
Elle communique avec l'API FastAPI via HTTP.

% =====================================================================
% 7. MONITORING
% =====================================================================
\section{Monitoring}

\subsection{Metriques exposees}

L'API expose 9 metriques Prometheus via \texttt{/metrics} :

\begin{table}[h]
\centering
\small
\begin{tabular}{lll}
\toprule
\textbf{Metrique} & \textbf{Type} & \textbf{Description} \\
\midrule
\texttt{prediction\_requests\_total} & Counter & Nombre total de requetes \\
\texttt{prediction\_errors\_total} & Counter & Nombre total d'erreurs \\
\texttt{prediction\_latency\_seconds} & Histogram & Latence des predictions \\
\texttt{model\_rmse\_price} & Gauge & RMSE du modele charge \\
\texttt{model\_r2\_price} & Gauge & R\textsuperscript{2} du modele \\
\texttt{model\_mae\_price} & Gauge & MAE du modele \\
\texttt{model\_mape} & Gauge & MAPE du modele \\
\texttt{training\_runs\_total} & Gauge & Nombre de runs d'entrainement \\
\texttt{data\_rows\_processed} & Gauge & Lignes de donnees traitees \\
\bottomrule
\end{tabular}
\caption{Metriques Prometheus exposees par l'API}
\end{table}

\subsection{Dashboard Grafana}

Le dashboard Grafana est structure en 4 sections :
\begin{enumerate}
    \item \textbf{KPIs Modele} -- RMSE, MAE, R\textsuperscript{2}, MAPE (vue d'ensemble)
    \item \textbf{Performance API} -- Compteurs de requetes/erreurs, taux de requetes
    \item \textbf{Latence} -- Percentiles p50, p95, p99
    \item \textbf{Entrainement} -- Nombre de runs, lignes de donnees traitees
\end{enumerate}

% =====================================================================
% 8. DEPLOIEMENT
% =====================================================================
\section{Deploiement}

\subsection{Docker Compose (developpement)}

L'environnement de dev comprend 8 services orchestres par Docker Compose :
PostgreSQL, MinIO, Airflow, FastAPI, Gradio, Prometheus, Grafana.

Les variables d'environnement sont centralisees via un ancrage YAML
(\texttt{x-common-env}) et les secrets sont lus depuis le fichier \texttt{.env} local.

\subsection{Kubernetes (production)}

Le deploiement K8s utilise :
\begin{itemize}
    \item \textbf{ConfigMaps} pour la configuration non-sensible (URIs, noms d'experiences)
    \item \textbf{Secrets} (\texttt{app-secrets}) pour les credentials (Postgres, MinIO, Kaggle, SMTP)
    \item \textbf{PersistentVolumeClaims} pour le stockage persistant (PostgreSQL, MinIO, MLflow, Airflow logs)
    \item \textbf{NodePort Services} pour l'exposition des services
    \item \textbf{Health probes} (readiness + liveness) sur chaque deployment
\end{itemize}

\subsection{Gestion des secrets}

\begin{table}[h]
\centering
\small
\begin{tabular}{lll}
\toprule
\textbf{Contexte} & \textbf{Mecanisme} & \textbf{Fichier} \\
\midrule
Dev local (sans Docker) & \texttt{.env} + \texttt{python-dotenv} & \texttt{.env} \\
Docker Compose & \texttt{.env} auto-charge par Compose & \texttt{.env} + \texttt{docker-compose.yml} \\
Kubernetes & ConfigMap + Secret & \texttt{k8s/secrets.yaml} \\
GitHub Actions & Repository Secrets & Settings $\rightarrow$ Secrets \\
\bottomrule
\end{tabular}
\caption{Gestion des secrets par environnement}
\end{table}

% =====================================================================
% 9. CI/CD
% =====================================================================
\section{CI/CD}

\subsection{Pipeline CI}

Declenchee sur chaque push/PR vers \texttt{main} ou \texttt{dev} :
\begin{enumerate}
    \item \textbf{Lint} -- Verification du code avec Ruff (linting + formatting)
    \item \textbf{Tests} -- Execution de pytest avec rapport de couverture
    \item \textbf{Docker Build} -- Verification que les images Docker se construisent
\end{enumerate}

\subsection{Pipeline CD}

Declenchee sur chaque tag \texttt{v*} :
\begin{enumerate}
    \item Build et push des images Docker vers GitHub Container Registry (GHCR)
    \item Deploiement sur Kubernetes via \texttt{kubectl apply}
    \item Verification du rollout (\texttt{kubectl rollout status})
\end{enumerate}

% =====================================================================
% 10. CONCLUSION
% =====================================================================
\section{Conclusion}

Ce projet demontre la mise en place d'une plateforme MLOps complete couvrant
l'ensemble du cycle de vie d'un modele de Machine Learning. Les choix
architecturaux (decoupage en microservices, separation des concerns, infrastructure
as code) permettent une evolution independante de chaque composant.

Le mecanisme de \textbf{continuous training avec comparaison de modeles} assure
que le modele en production ne se degrade jamais, garantissant la fiabilite du
systeme en conditions reelles.

Les axes d'amelioration possibles incluent :
\begin{itemize}
    \item A/B testing entre modeles en production
    \item Monitoring de data drift (Evidently, Great Expectations)
    \item Autoscaling Kubernetes (HPA) base sur la latence
    \item Feature store pour la reutilisation des features
\end{itemize}

\end{document}
